% !TeX root = ../main.tex
\chapter{رگرسیون؛ پیش‌بینی پاسخ عددی}
\source{اسلاید رگرسیون، PDF ۳ تا ۵۵؛ بسط مفاهیم تعمیم و اریبی–واریانس از بخش‌های ۳٫۴ تا ۳٫۸ و ۶٫۱۰٫۳ کتاب. توضیح ridge ، lasso ، ایزوتونیک و رگرسیون محلی بر اسلاید و استدلال ریاضی تکیه دارد.}
\section{مسئله چیست؟}

در طبقه‌بندی می‌پرسیم «این خانه گران است یا ارزان؟» و در رگرسیون می‌پرسیم «قیمت این خانه چقدر است؟». دادهٔ آموزش زوج‌های \((\boldsymbol{x}_i,y_i)\) دارد؛ بردار \(\boldsymbol{x}_i\) شامل ویژگی‌هایی مانند مساحت، تعداد اتاق و فاصله تا حمل‌ونقل و \(y_i\) قیمت مشاهده‌شده است. خروجی مدل \(\widehat y_i=f(\boldsymbol{x}_i)\) پیش‌بینی است و با مقدار واقعی یکی نیست. بهتر است در جزوه از \(n\) برای تعداد نمونه و \(p\) برای تعداد ویژگی استفاده شود تا نمادها قاطی نشوند.

مدل با کمینه‌کردن خطا روی دادهٔ آموزش ساخته می‌شود، ولی هدف نهایی کم‌بودن خطا روی خانه‌ها یا شرایطی است که ندیده‌ایم. جداکردن دادهٔ آزمون باید پیش از مقیاس‌بندی، حذف برون‌افتاده، انتخاب ویژگی و انتخاب پارامتر انجام شود. آمارِ مقیاس‌بندی فقط از آموزش به دست می‌آید. برای دادهٔ زمانی، تقسیم زمانی به‌جای تقسیم تصادفی می‌تواند ضروری باشد. این ادامهٔ اصول بخش‌های ۳٫۵--۳٫۸ کتاب است.

\section{رگرسیون نزدیک‌ترین همسایه}

برای برآورد دمای محلی که ایستگاه ندارد، دمای ایستگاه‌های نزدیک را میانگین می‌گیریم. اگر پنج دمای نزدیک ۱۸، ۲۰، ۲۱، ۲۲ و ۲۱ باشند، پیش‌بینی \(102/5=20.4\) درجه است؛ این همان مثال اسلاید ۸ است. در حالت کلی:

\[
\widehat y(\boldsymbol{x})=\frac1k\sum_{i\in N_k(\boldsymbol{x})}y_i,
\qquad
\widehat y_w(\boldsymbol{x})=\frac{\sum_{i\in N_k}w_i y_i}{\sum_{i\in N_k}w_i},
\quad w_i=\frac1{d(\boldsymbol{x},\boldsymbol{x}_i)+\varepsilon}.
\]

ابتدا فاصله را تعریف می‌کنیم، سپس k همسایه را پیدا می‌کنیم و در پایان مقدار عددی آنان را ترکیب می‌کنیم. برای فاصلهٔ صفر باید صریحاً خروجی نمونه‌های هم‌مکان را استفاده کرد یا قرارداد پایدار داشت. k کوچک به جزئیات و نویز حساس‌تر است؛ k بزرگ نمودار را هموارتر می‌کند و ممکن است الگوی محلی را محو کند. مقیاس ویژگی‌ها مستقیماً روی انتخاب همسایه اثر می‌گذارد: «حقوق به ریال» می‌تواند «سن به سال» را تحت‌الشعاع قرار دهد. k را روی اعتبارسنجی انتخاب می‌کنیم. در میانگین با وزن‌های نامنفی، پیش‌بینی در بازهٔ مقادیر هدفِ همسایگان می‌ماند؛ به همین دلیل مدل معمولاً برون‌یابی عددی نمی‌کند.

\section{معیارهای خطا و دام همبستگی}

اگر \(e_i=y_i-\widehat y_i\)، آنگاه

\[
\operatorname{MAE}=n^{-1}\sum_i|e_i|,\quad
\operatorname{MSE}=n^{-1}\sum_i e_i^2,\quad
\operatorname{RMSE}=\sqrt{\operatorname{MSE}}.
\]

برای واقعی‌های \((2,4,6)\) و پیش‌بینی‌های \((3,4,4)\)، خطاها \((-1,0,2)\)، MAE برابر ۱، MSE برابر \(5/3\) و RMSE حدود ۱٫۲۹۱ است. MAE و RMSE واحد هدف را دارند؛ MSE واحدِ مربع هدف را دارد. مربع‌کردن خطا سبب می‌شود یک خطای بزرگ نسبت به چند خطای کوچک وزن بیشتری بگیرد. به همین علت گزارش هم‌زمان MAE و RMSE می‌تواند وجود خطاهای بزرگ را آشکار کند.

همبستگی پیرسون \[
r=\frac{\sum_i(y_i-\bar y)(\widehat y_i-\overline{\widehat y})}{\sqrt{\sum_i(y_i-\bar y)^2\sum_i(\widehat y_i-\overline{\widehat y})^2}}
\] روند خطی را می‌سنجد، نه توافق عددی را. برای واقعی‌های ۱،۲،۳ و پیش‌بینی‌های ۱۰۰۰،۲۰۰۰،۳۰۰۰، همبستگی ۱ است اما خطا بسیار بزرگ است. اگر یکی از دو بردار ثابت باشد، مخرج صفر و همبستگی تعریف‌نشده است.

تعریف عمومی ضریب تعیین چنین است: \[
R^2=1-\frac{\sum_i(y_i-\widehat y_i)^2}{\sum_i(y_i-\bar y)^2}.
\] برای مثال بالا \(\bar y=4\)، SST برابر ۸ و SSE برابر ۵ است، پس \(R^2=0.375\). پیش‌بینی ثابتِ میانگین دادهٔ ارزیابی $R^2$ صفر دارد و پیش‌بینی کاملاً درست $R^2$ یک دارد. \textbf{اصلاح اسلاید ۱۲:} $R^2$ به‌طور عمومی محدود به بازهٔ صفر و یک نیست؛ برای مدل بد روی آزمون منفی می‌شود. نسبت explained SS / total SS فقط در شرایط تجزیهٔ مناسب، از جمله OLS دارای عرض از مبدأ روی آموزش، با تعریف عمومی برابر است. وقتی تمام y ها یکسان‌اند، SST صفر است و باید قرارداد نرم‌افزاری را جداگانه توضیح داد.

\section{درخت رگرسیون و درخت مدل}

ساختار سؤال‌های درخت تغییر نمی‌کند: مثلاً «مساحت کمتر از ۱۰۰ است؟». تفاوت در هدف شکستن و مقدار برگ است. در برگ با کمینه‌سازی SSE باید میانگین y های آموزش را پیش‌بینی کنیم. برای هر شکستنِ کاندید، کاهش خطا را می‌سنجیم: \[
\Delta=\operatorname{SSE}(P)-\operatorname{SSE}(L)-\operatorname{SSE}(R).
\]

مثال مستقل: \(x=(1,2,3,4)\)، \(y=(1,2,8,9)\). پیش از شکستن، میانگین ۵ و SSE برابر \(16+9+9+16=50\) است. شکستن در ۲٫۵ دو میانگین ۱٫۵ و ۸٫۵ می‌دهد و SSE کل ۱ می‌شود؛ کاهش برابر ۴۹ است. نمودار خروجی پله‌ای یا تکه‌ای ثابت خواهد بود. عمق زیاد و برگ‌های تک‌نمونه‌ای می‌توانند نویز را حفظ کنند؛ محدودیت عمق و حداقل اندازهٔ برگ ابزار کنترل‌اند.

در «درخت مدل» به‌جای ثابت در هر برگ، یک رگرسیون خطی محلی آموزش می‌دهیم؛ مثلاً مصرف خودروهای دیزلی و بنزینی هرکدام رابطهٔ متفاوتی با وزن دارند. هنگام پیش‌بینی ابتدا مسیر درخت و سپس فرمول برگ اجرا می‌شود. چنین مدلی تکه‌ای خطی است، و برخلاف برگ‌های میانگینِ درخت رگرسیون، ممکن است برون‌یابی کند.

\section{رگرسیون خطی و کمترین مربعات}

مدل ساده \(\widehat y=w_0+w_1x\) و مدل چندمتغیره \(\widehat y=w_0+\sum_{j=1}^p w_jx_j\) است. «خطی» در اینجا یعنی خطی نسبت به ضرایب؛ می‌توان ویژگیِ \(x^2\) را هم وارد مدل خطی کرد. تابع هدف OLS برابر \(\min_{\boldsymbol w}\|\boldsymbol y-X\boldsymbol w\|_2^2\) است؛ ستون ثابت در X عرض از مبدأ را نمایش می‌دهد. در حالت ساده: \[
\widehat w_1=\frac{\sum_i(x_i-\bar x)(y_i-\bar y)}{\sum_i(x_i-\bar x)^2},
\qquad \widehat w_0=\bar y-\widehat w_1\bar x.
\]

برای \(x=(1,2,3),y=(2,3,5)\)، \(\bar x=2\)، \(\bar y=10/3\)، شیب \(3/2\) و عرض \(1/3\) است. پیش‌بینی‌ها \((11/6,10/3,29/6)\) و SSE برابر \(1/6\) می‌شوند. در چندمتغیره اگر X رتبهٔ ستونی کامل داشته باشد، راه‌حل یکتا و به‌صورت \((X^TX)^{-1}X^Ty\) بیان می‌شود؛ در اجرا QR/SVD پایدارتر از معکوس‌گیری صریح است.

برای برازش عددیِ کمترین مربعات، نرمال‌بودن خطا لازم نیست. برای تفسیر آماری کلاسیک، درست‌بودن شکل میانگین شرطی، صفر بودن میانگین خطای شرطی، استقلال مناسب مشاهدات و شرایط واریانس اهمیت دارند؛ نرمال‌بودن معمولاً به استنباط دقیق کوچک‌نمونه مربوط است. همبستگی شدید ویژگی‌ها می‌تواند ضرایب را ناپایدار کند. کوچک‌بودن SSE آموزش به‌تنهایی دلیل تعمیم خوب نیست.

\section{منظم‌سازی ridge و lasso}

افزودن ویژگی‌ها انعطاف مدل را زیاد می‌کند. منظم‌سازی در مقابل بهبود برازش، برای ضرایب بزرگ هزینه تعیین می‌کند: \[
\text{Ridge}:\ \min_w \sum_i(y_i-\widehat y_i)^2+\lambda\sum_{j=1}^p w_j^2,
\qquad
\text{Lasso}:\ \min_w \sum_i(y_i-\widehat y_i)^2+\lambda\sum_{j=1}^p|w_j|.
\]

عرض از مبدأ معمولاً جریمه نمی‌شود. مقیاس‌بندی ویژگی‌ها پیش از جریمه مهم است؛ در غیر این صورت اندازهٔ ضریب به واحد اندازه‌گیری وابسته می‌شود. Ridge ضرایب را پیوسته کوچک می‌کند؛ lasso می‌تواند ضرایب را دقیقاً صفر و نوعی انتخاب ویژگی انجام دهد. در حضور ویژگی‌های بسیار هم‌بسته، انتخاب lasso ممکن است ناپایدار باشد. \(\lambda=0\) حالت بدون جریمه است و الزاماً بیش‌برازش ایجاد نمی‌کند؛ \(\lambda\) خیلی بزرگ، با عرض از مبدأ آزاد، مدل را به پیش‌بینی نزدیک میانگین می‌رساند. «بدون جریمه» با «بدون نرمال‌سازی» یکسان نیست؛ تعبیر اسلاید ۲۱ باید دقیق‌تر نوشته شود.

برای شهود، مسئلهٔ تک‌ضریبی \((w-3)^2+2w^2\) کمینهٔ ridge برابر ۱ دارد؛ اما \((w-3)^2+2|w|\) کمینهٔ lasso برابر ۲ دارد. در قرارداد کلی \((w-a)^2+\lambda|w|\)، پاسخ \(\operatorname{sign}(a)\max(|a|-\lambda/2,0)\) است. ضریب \(1/2\) به قرارداد تابع هدف وابسته است؛ در کتابخانه‌ها ممکن است ضرایب n یا ۱/۲ متفاوت باشند.

\section{درون‌یابی، برون‌یابی و خط مبنا}

در کاربرد علمی، درون‌یابی/برون‌یابی معمولاً نسبت به ناحیهٔ ورودی آموزش تعریف می‌شود. اسلاید ۲۳ علاوه بر آن روی خارج‌شدن مقدار پیش‌بینی از بازهٔ هدف تأکید دارد؛ این دو مفهوم را نباید یکی گرفت. خطی می‌تواند بیرون بازهٔ y آموزش پیش‌بینی کند؛ kNN ِ میانگین و درختِ میانگین نمی‌توانند، اما این ویژگی تضمین نمی‌کند پیش‌بینی آنان در موقعیت جدید معتبر باشد. حتی در بازهٔ ورودی هم نقاط کم‌تراکم ممکن است نامطمئن باشند. برون‌یابی برای آینده‌نگری لازم است ولی به دوام فرض رابطه تکیه دارد.

خط مبنا باید پیش از مدل پیچیده ساخته شود. ثابت میانگین آموزش SSE را و ثابت میانه MAE را روی آموزش کمینه می‌کند. در آزمون فقط همان ثابت آموخته‌شده از آموزش به کار می‌رود، نه میانگینِ هدف آزمون. اگر مدل از این مبنا بهتر نباشد، پیچیدگی آن توجیه تجربی ندارد.

\section{ایزوتونیک و الگوریتم ادغام ناقضان مجاور}

گاهی رابطه خطی نیست ولی می‌دانیم افزایشی است؛ مثلاً انتظار داریم یک امتیاز کالیبره‌شده با امتیاز خام کاهش نیابد. پس پس از مرتب‌کردن x ها، حل می‌کنیم: \[
\min_{z_1\le\cdots\le z_n}\sum_{i=1}^n a_i(y_i-z_i)^2,\qquad a_i>0.
\] هر مشاهده را ابتدا یک بلوک می‌گیریم. اگر مقدار بلوک چپ از بلوک راست بیشتر باشد، دو بلوک ناقض‌اند؛ آن‌ها را ادغام و میانگین وزنی‌شان را جایگزین می‌کنیم. سپس دوباره بلوک قبل را بررسی می‌کنیم و تا رفع همهٔ نقض‌ها ادامه می‌دهیم. وزن بلوک ادغام‌شده مجموع وزن‌هاست؛ فراموش‌کردن وزن سبب خطا می‌شود.

مثال: y های مرتب \((1,4,3,2,5)\). ابتدا ۴ و ۳ به میانگین ۳٫۵ با وزن ۲ تبدیل می‌شوند. چون ۳٫۵ از ۲ بزرگ‌تر است، بلوک سه‌تایی میانگین \((4+3+2)/3=3\) می‌گیرد. جواب \((1,3,3,3,5)\) و SSE برابر ۲ است. نقاط x جابه‌جا نمی‌شوند؛ مقادیر برازش تعدیل می‌شوند. برابری برای همهٔ همسایه‌ها لازم نیست؛ شرط صحیح \(z_i\le z_{i+1}\) است. نمایش بین گره‌ها می‌تواند با قرارداد خطی یا پله‌ای انجام شود. اسلاید ۳۶ رفتار بیرون دامنه را با clip/nan/raise تفکیک می‌کند؛ این رفتار بخشی از تعریف کاربردی مدل است.

\section{چندجمله‌ای و رگرسیون محلی}

مدل درجه d برابر \(\widehat y=\sum_{j=0}^d w_jx^j\) است و همچنان نسبت به w خطی است. افزایش d امکان دنبال‌کردن خمیدگی و همچنین نویز را فراهم می‌کند. برای p ویژگی، تعداد تک‌جمله‌ای‌های تا درجه d با جملهٔ ثابت \(\binom{p+d}{d}\) است؛ برای $p=3$ و $d=3$ تعداد ۲۰ داریم: یک ثابت، سه درجه‌یک، شش درجه‌دو و ده درجه‌سه. این شمارش، توضیح دقیق‌تر اسلاید ۳۹ است. داشتن ویژگی بیشتر از نمونه به‌خودی‌خود ممنوعیت ریاضی نیست، ولی بدون منظم‌سازی/ساختار مناسب، ناپایداری و بیش‌برازش محتمل می‌شود. انتخاب درجه باید با اعتبارسنجی انجام شود. بیرون دامنهٔ آموزش نوسان چندجمله‌ای می‌تواند شدید شود.

رگرسیون محلی برای هر x تازه همسایه‌های نزدیک را پیدا می‌کند، روی همان همسایه‌ها خط برازش می‌دهد و در x تازه آن خط را ارزیابی می‌کند. شکل عمومی کمترین مربعات وزن‌دار است: \[
\widehat\beta(x)=\arg\min_\beta\sum_i K\!\left(\frac{x_i-x}{h}\right)(y_i-\beta_0-\beta_1(x_i-x))^2,
\qquad\widehat y(x)=\widehat\beta_0(x).
\] نسخهٔ اسلاید فقط k همسایه را به کار می‌گیرد. h یا k میزان هموارسازی را کنترل می‌کند. برخلاف kNN ، خط محلی روند را هم تخمین می‌زند و ممکن است از بازهٔ هدف همسایه‌ها خارج شود. فرض شهودی این است که تابعِ خمیده در همسایگی کوچک تقریباً خطی است. هزینهٔ پیش‌بینی بیشتر است، چون بازیابی همسایه و برازش برای هر پرسش تکرار می‌شود.

\section{شبکهٔ عصبی و موازنهٔ اریبی/واریانس}

برای تبدیل شبکهٔ طبقه‌بندی به رگرسیون، خروجی عددی را نگه می‌داریم و به‌جای خطای طبقه‌بندی، مثلاً MSE را بهینه می‌کنیم. لایهٔ خروجی خطی محدود به بازهٔ صفر و یک نیست؛ عبارت اسلاید ۴۵ این محدودیت را توجیه نمی‌کند. شبکهٔ چندلایه تنها وقتی از ترکیب خطی فراتر می‌رود که میان لایه‌ها فعال‌ساز غیرخطی داشته باشیم. افزودن صرفِ لایه‌های خطی دوباره یک تابع خطی می‌سازد.

برای مدل داده \(Y=f(x)+\varepsilon\)، با \(E[\varepsilon\mid x]=0\)، و آموزش تصادفی D ، تجزیهٔ خطای مربعی در x چنین است: \[
E_{D,\varepsilon}[(Y-\widehat f_D(x))^2]
=\underbrace{(E_D\widehat f_D(x)-f(x))^2}_{\mathrm{Bias}^2}
+\underbrace{E_D[(\widehat f_D(x)-E_D\widehat f_D(x))^2]}_{\mathrm{Variance}}
+\underbrace{\operatorname{Var}(\varepsilon\mid x)}_{\mathrm{Noise}}.
\] اریبی یعنی میانگین مدل‌های آموزش‌دیده از رابطهٔ واقعی فاصله دارد؛ واریانس یعنی مدل با عوض‌شدن نمونهٔ آموزش تغییر زیادی می‌کند. نویزِ کاهش‌ناپذیر تحت اطلاعات و مدل دادهٔ فعلی باقی می‌ماند؛ حذف کورکورانهٔ نقاط غیرعادی آن را لزوماً از بین نمی‌برد. مثال حسگر خراب می‌تواند آلودگی قابل‌اصلاح باشد و با تصادفی‌بودن ذاتی یکی نیست.

منحنی یادگیری خطای آموزش و اعتبارسنجی را در برابر اندازهٔ آموزش می‌کشد. خطای هر دو زیاد و نزدیک، نشانهٔ احتمالی کم‌برازش است؛ خطای آموزش کم ولی فاصلهٔ بزرگ با اعتبارسنجی، نشانهٔ احتمالی بیش‌برازش است. دادهٔ بیشتر اغلب واریانس را کم می‌کند اما شکل تابع اشتباه را اصلاح نمی‌کند. اسلاید ۵۲ امکان آستانه‌گذاری خروجی رگرسیون روی ۰٫۵ برای طبقه‌بندی دودویی را مطرح می‌کند؛ خروجی رگرسیون خطی الزاماً احتمال معتبر نیست و برچسب‌گذاری عددی کلاس‌های اسمی فاصلهٔ مصنوعی ایجاد می‌کند.
\section{چرا میانگین، پاسخِ کمترین مربعات است؟}
در یک برگ درخت، یک بلوک ایزوتونیک یا مدل ثابت، می‌خواهیم یک عدد $c$ را به همهٔ پاسخ‌ها نسبت دهیم. با زیان مربعی:
\[L(c)=\sum_{i=1}^n(y_i-c)^2,\qquad L'(c)=2nc-2\sum_i y_i.\]
مشتق صفر می‌دهد $c=\bar y$ و مشتق دوم $2n>0$ کمینه بودن را ثابت می‌کند. با وزن‌های مثبت، همین استدلال میانگین وزن‌دار را می‌دهد. اما برای زیان قدرمطلق، یک میانه کمینه‌کننده است. انتخاب نماینده از تابع زیان می‌آید، نه از یک قاعدهٔ ثابت برای تمام الگوریتم‌ها.

برای پاسخ‌های $1,2,9$، بهترین ثابت مربعی ۴ و مجموع مربع خطا $9+4+25=38$ است. میانه ۲، مجموع قدرمطلق خطای $1+0+7=8$ می‌دهد؛ در مقابل، مجموع قدرمطلق خطای ثابت ۴ برابر ۱۰ است. عوض کردن معیار می‌تواند بهترین مدل را هم عوض کند. وقتی RMSE و MAE اختلاف زیادی دارند، چند باقیماندهٔ بزرگ را بررسی کنید؛ حذف آن‌ها بدون بررسی منشأ، پاسخ علمی نیست.

\subsection*{ارزیابی یک برازش، فراتر از رسم منحنی}
باقیمانده را در برابر ورودی و پیش‌بینی رسم می‌کنیم. الگوی خمیده می‌تواند نشان دهد رابطهٔ خطی کافی نیست؛ افزایش دامنهٔ پراکندگی می‌تواند نشانهٔ واریانس ناهمسان باشد؛ ترتیب زمانی باقیمانده‌ها می‌تواند وابستگیِ نادیده را آشکار کند. این‌ها ابزار تشخیص‌اند و به‌تنهایی اثبات یک مدل خاص نیستند. پیش‌بینی، تفسیر ضریب و استنباط آماری سه هدف متفاوت‌اند؛ برای نتیجه‌گیری احتمالی دربارهٔ ضریب به فرض‌های خطا و طراحی داده نیز نیاز داریم.

در مقایسهٔ درجه‌ها، یک مجموعهٔ آموزش و یک اعتبارسنجی ثابت نگه دارید. اگر دادهٔ تصادفی با هر حرکت کنترل تازه شود، معلوم نیست تغییر امتیاز از درجه آمده یا از عوض شدن نمونه. در آزمایشگاه این فصل بذر ثابت است، جریمه روی عرض از مبدأ اعمال نمی‌شود و قرارداد تابع هدف، مجموع مربع خطا به‌علاوهٔ جریمه است. بنابراین مقدار عددی جریمه را بدون توجه به قرارداد با نرم‌افزار دیگر مقایسه نکنید.

\section{خودآزمایی و تمرین}
\begin{selftests}
\begin{selftest}{خطاها و واحدها را محاسبه کنید}
برای حقیقت $(2,4,6)$ و پیش‌بینی $(3,4,4)$، MAE، MSE و RMSE چیست؟
\answer{خطاها $(-1,0,2)$، مقدار MAE برابر یک، MSE برابر $5/3$ و RMSE برابر $\sqrt{5/3}\approx1.2910$ است. MAE و RMSE واحد پاسخ و MSE مربع واحد پاسخ را دارد.}
\end{selftest}
\begin{selftest}{همبستگی کامل، پیش‌بینی کامل؟}
اگر $\widehat y=2y+10$ باشد، چرا ممکن است همبستگی یک ولی خطا زیاد باشد؟
\answer{همبستگی ترتیب و رابطهٔ خطی را می‌سنجد و به انتقال و مقیاس مثبت حساس نیست؛ ولی اختلاف با مقدار واقعی بزرگ است. برای کالیبراسیون عددی باید خطا و خط مبنا را هم بررسی کرد.}
\end{selftest}
\begin{selftest}{R² منفی را تفسیر کنید}
اگر $SSE=20$ و $SST=8$ باشد، $R^2$ چیست؟
\answer{$R^2=1-20/8=-1.5$ است؛ مدل با این تعریف از پیش‌بینی ثابتِ میانگین همان دادهٔ ارزیابی بدتر است. اگر پاسخ‌ها ثابت باشند، SST صفر و این تعریفِ کسرشکل نامعتبر است.}
\end{selftest}
\begin{selftest}{میانگین یا میانه؟}
برای $1,2,9$ کدام ثابت برای مربع خطا و کدام برای قدرمطلق خطا بهینه است؟
\answer{برای مربع خطا، میانگین ۴ و برای قدرمطلق خطا، میانه ۲ است. تابع زیان تعیین می‌کند «بهترین ثابت» چه معنایی دارد.}
\end{selftest}
\begin{selftest}{ادغام ایزوتونیک را انجام دهید}
برازش غیرکاهشیِ $1,4,3,2,5$ با وزن برابر چیست؟
\answer{۴ و ۳ ابتدا بلوک $3.5$ می‌سازند؛ این بلوک با ۲ ناسازگار است و سه مقدار به میانگین ۳ ادغام می‌شوند. پاسخ $(1,3,3,3,5)$ و مجموع مربع خطا ۲ است؛ ترتیب ورودی‌ها عوض نمی‌شود.}
\end{selftest}
\begin{selftest}{درجه، جریمه و برون‌یابی را نقد کنید}
درجهٔ هفت روی آموزش خطای کمتری دارد ولی بیرون بازه نوسان می‌کند. آیا باید آن را انتخاب کنیم؟
\answer{خطای آموزش کافی نیست. درجه و جریمه را با اعتبارسنجی و دامنهٔ کاربرد انتخاب می‌کنیم؛ عملکرد خوب داخل بازه تضمین رفتار خارج آن نیست. استانداردسازی و ساخت ویژگی هم باید درون خط لولهٔ آموزش باشد.}
\end{selftest}
\end{selftests}
